\section{Poisson and Hawkes}
\label{sec:poisson-hawkes}

We need two building blocks: the Poisson process (which the MBPP will be) and the
Hawkes process (which we want to fit). The crucial contrast is whether the
\emph{compensator} is deterministic.

\subsection{Counting processes, intensity, compensator}

A simple temporal point process is described by its counting process
$N(t)=\#\{i:t_i\le t\}$, adapted to its internal history
$\Hist_t=\sigma\{N(s):s\le t\}$. By the Doob--Meyer decomposition there is a unique
predictable increasing \emph{compensator} $\Lambda$ with
\begin{equation}
M(t)=N(t)-\Lambda(t)\ \text{an }(\Hist_t)\text{-local martingale.}
\end{equation}
When $\Lambda$ is absolutely continuous, $\Lambda(t)=\int_0^t\lambda$, and the
density $\lambda$ is the \emph{conditional intensity}
$\lambda(t)=\lim_{h\downarrow0}\frac1h\,\E[N(t{+}h)-N(t)\mid\Hist_{t^-}]$.

\begin{intuition}
$\Lambda$ is the predictable ``running forecast'' of how many events should have
happened by time $t$; $\lambda$ is that forecast's instantaneous rate. Subtracting
the forecast from reality leaves pure surprise --- the martingale $M$ --- so
$\lambda$ captures exactly the part of the future the past lets us anticipate.
\end{intuition}

\subsection{The Poisson process and why deterministic $\Lambda$ matters}

\begin{definition}[Poisson process]
$N$ is an inhomogeneous Poisson process with mean measure $\Lambda$ if for disjoint
bounded sets $A_1,\dots,A_k$ the counts $N(A_i)$ are independent with
$N(A_i)\sim\mathrm{Poisson}(\Lambda(A_i))$.
\end{definition}

Two consequences are used constantly: a Poisson process has \emph{independent
increments}, and $\E[N(t)]=\Lambda(t)$. The deep fact we rely on is the converse.

\begin{theorem}[Watanabe characterisation]\label{thm:watanabe}
Let $N$ be a simple point process with continuous compensator $\Lambda$. Then $N$
is an inhomogeneous Poisson process with mean measure $\Lambda$ \emph{if and only
if $\Lambda$ is deterministic.}
\end{theorem}

\begin{intuition}
Dependence between increments can only come from the intensity \emph{reacting to
history}. If the intensity --- hence the compensator --- is fixed in advance, there
is no channel through which the past can speak to the future, so the process must
be Poisson. \textbf{This single fact is why the MBPP works}: replacing the random
Hawkes intensity $\lambda$ by its deterministic mean $\xi$ severs the feedback loop
and gives us back independent increments (and therefore a product likelihood,
\S\ref{sec:no-covariates}).
\end{intuition}

\subsection{The likelihood of a point process}

\begin{theorem}[Jacod formula]\label{thm:jacod}
A simple point process with $(\Prob,\Hist_t)$-intensity $\lambda$ has, on $[0,T]$,
the log-likelihood (relative to a unit-rate Poisson reference, up to a constant)
\begin{equation}
\ell(\theta)=\sum_{t_i\le T}\log\lambda(t_i;\theta)-\int_0^T\lambda(s;\theta)\dd s .
\label{eq:ppll}
\end{equation}
\end{theorem}

\begin{intuition}
The likelihood rewards forecasting high intensity exactly where events landed
($\sum\log\lambda(t_i)$) and penalises total forecast mass ($\int\lambda$): ``put
probability on the bumps, don't waste it elsewhere.'' It is the point-process
analogue of $\prod_i p(x_i)$, and it manifestly needs the times $t_i$ --- the
obstruction in the interval-censored setting.
\end{intuition}

\subsection{The Hawkes process and why it is \emph{not} Poisson}

\begin{model}[Hawkes process]\label{mod:hawkes}
The Hawkes process has stochastic intensity
\begin{equation}
\lambda(t)=s(t)+\int_0^{t^-}\!\phi(t-u)\dd N(u)
=s(t)+\sum_{t_i<t}\phi(t-t_i),
\label{eq:hawkes}
\end{equation}
with exogenous (immigrant) rate $s$ and triggering kernel $\phi$. The
\emph{exponential kernel} is $\phi(t)=\kappa\theta\,e^{-\theta t}\,\one[t>0]$, for
which $\int_0^\infty\phi=\kappa$ is the \emph{branching ratio} and $\theta$ the
decay rate.
\end{model}

Because $\lambda$ depends on the realised history, the compensator
$\Lambda(t)=\int_0^t\lambda$ is \emph{random}; by Theorem~\ref{thm:watanabe} the
Hawkes process is therefore \emph{not} Poisson and has dependent increments.

\subsection{The cluster representation and the mean-intensity equation}

\begin{theorem}[Hawkes--Oakes cluster representation]\label{thm:cluster}
The Hawkes process of \eqref{eq:hawkes} is distributionally equal to the following
branching construction: \emph{(i)} immigrants arrive as a Poisson process of
intensity $s$; \emph{(ii)} each event at time $u$ independently produces direct
offspring as a Poisson process of intensity $\phi(\cdot-u)$ on $(u,\infty)$;
\emph{(iii)} the process is the superposition of all immigrants and descendants.
\end{theorem}

\begin{intuition}
Self-excitation made literal: every event founds a ``family,'' and the whole
process is a \emph{forest of independent family trees} rooted at immigrants. The
mean number of direct children per event is $\int\phi=\kappa$, so a family is a
Galton--Watson tree with mean offspring $\kappa$ --- finite on average iff
$\kappa<1$ (\emph{subcritical}), with expected total size $1/(1-\kappa)$.
\end{intuition}

Averaging the cluster representation gives the equation at the heart of everything.

\begin{proposition}[Mean-intensity / MBPP equation]\label{prop:mbpp-eq}
Under the standing assumptions of \S\ref{sec:assumptions}, the expected intensity
$\xi(t)=\E[\lambda(t)]$ satisfies
\begin{equation}
\xi(t)=s(t)+\int_0^t\phi(t-u)\,\xi(u)\dd u .
\label{eq:mbpp}
\end{equation}
\end{proposition}
\begin{proof}
Take $\E[\cdot]$ in \eqref{eq:hawkes}. By Tonelli (non-negative integrands) the
expectation passes inside, and using the tower property
$\E[\dd N(u)]=\E[\E[\dd N(u)\mid\Hist_{u^-}]]=\E[\lambda(u)]\dd u=\xi(u)\dd u$.
\end{proof}

\begin{intuition}
Average the \emph{stochastic} feedback over all realisations and it becomes
\emph{deterministic} feedback: ``on average, today's rate is the baseline plus the
kernel-weighted echo of the average past.'' Equation~\eqref{eq:mbpp} is the object
we will solve again and again. It is a \emph{Volterra integral equation of the
second kind} (Appendix~\ref{app:volterra}); when --- as here --- the kernel depends
only on $t-u$, it is moreover a \emph{convolution}, which is what makes the next
section entirely closed-form.
\end{intuition}

\begin{remark}[Why interval censoring is hard for Hawkes directly]
Decomposing a bucket count over offspring generations produces nested
\emph{Borel} convolutions with no finite closed form at finite $T$, and Monte
Carlo over the generation structure scales poorly. This is the formal sense in
which ``the Hawkes process cannot be fit interval-censored'' --- and the reason we
detour through the MBPP.
\end{remark}
